\documentclass[../../../include/open-logic-section]{subfiles}

\begin{document}

\olfileid[id]{sth}{card-arithmetic}{ch}

\olsection{Hipotesis Kontinuum}

Hasil sebelumnya mengisyaratkan (secara tepat) bahwa eksponensiasi kardinal
akan cukup \emph{mudah} jika kardinal tak hingga dijamin ``berinteraksi
secara langsung'' dengan pangkat $2$, yakni (berdasarkan
\olref[opps]{lem:SizePowerset2Exp}) dengan operasi mengambil himpunan kuasa.
Namun, kita tidak dapat mengandaikan bahwa kardinal tak hingga \emph{memang}
berinteraksi secara langsung dengan himpunan kuasa.

Untuk mulai menguraikan persoalan ini, kita memperkenalkan suatu notasi yang
berguna.

\begin{defn}
Jika $\cardsucc{\cardfont{a}}$ merupakan kardinal terkecil yang secara ketat
lebih besar daripada $\cardfont{a}$, kita mendefinisikan dua barisan tak
hingga:
\begin{align*}
	\aleph_{0} &\defis \omega & 		
	\beth_{0} &\defis \omega\\
	\aleph_{\alpha \ordplus 1} &\defis \cardsucc{(\aleph_{\alpha})} &
	\beth_{\alpha+1} &\defis \cardexpo{2}{\beth_{\alpha}}\\
	\aleph_{\alpha} &\defis \bigcup_{\beta< \alpha} \aleph_{\beta} &
	\beth_{\alpha} &\defis \bigcup_{\beta < \alpha}\beth_{\beta} & \text{jika $\alpha$ merupakan ordinal limit}.
	\end{align*}
\end{defn}

Definisi $\cardsucc{\cardfont{a}}$ terbentuk dengan baik, sebab
\olref[cardinals][classing]{lem:NoLargestCardinal} memberi tahu kita bahwa,
untuk setiap kardinal $\cardfont{a}$, terdapat suatu kardinal yang lebih
besar daripada $\cardfont{a}$, dan Induksi Transfinit menjamin bahwa terdapat
suatu kardinal \emph{terkecil} yang lebih besar daripada $\cardfont{a}$.
Sisa definisi kedua barisan tersebut diberikan melalui rekursi transfinit.

Cantor memperkenalkan notasi ``$\aleph$'' ini; lambang tersebut ialah
\emph{aleph}, huruf pertama dalam alfabet Ibrani dan huruf pertama dalam kata
Ibrani untuk ``tak hingga''. Peirce memperkenalkan notasi ``$\beth$'';
lambang tersebut ialah \emph{beth}, huruf kedua dalam alfabet
Ibrani.\footnote{Peirce menggunakan notasi ini dalam sepucuk surat kepada
Cantor pada Desember 1900. Sayangnya, dalam surat tersebut Peirce juga
memberikan argumen yang keliru bahwa $\beth_\alpha$ tidak ada untuk
$\alpha \geq \omega$.} Notasi-notasi ini memberi kita kardinal tak hingga.

\begin{prop}
$\aleph_\alpha$ dan $\beth_\alpha$ merupakan kardinal, untuk setiap ordinal
$\alpha$.
\end{prop}

\begin{proof}
Kedua hasil berlaku berdasarkan suatu induksi transfinit sederhana.
$\aleph_0 = \beth_0 = \omega$ merupakan suatu kardinal berdasarkan
\olref[cardinals][classing]{omegaisacardinal}. Dengan mengandaikan bahwa
$\aleph_\alpha$ dan $\beth_\alpha$ keduanya merupakan kardinal,
$\aleph_{\alpha+1}$ dan $\beth_{\alpha+1}$ secara eksplisit didefinisikan
sebagai kardinal. Selain itu, gabungan suatu himpunan kardinal merupakan suatu
kardinal berdasarkan
\olref[cardinals][classing]{unioncardinalscardinal}.
\end{proof}
\noindent
Lebih jauh, setiap kardinal tak hingga merupakan suatu $\aleph$:

\begin{prop}
Jika $\cardfont{a}$ suatu kardinal tak hingga, maka
$\cardfont{a} = \aleph_\gamma$ bagi suatu $\gamma$ yang unik.
\end{prop}

\begin{proof}
Berdasarkan induksi transfinit atas kardinal. Untuk induksi, andaikan bahwa
jika $\cardfont{b} < \cardfont{a}$, maka
$\cardfont{b} = \aleph_{\gamma_\cardfont{b}}$. Jika
$\cardfont{a} = \cardsucc{\cardfont{b}}$ bagi suatu $\cardfont{b}$, maka
$\cardfont{a} = \cardsucc{(\aleph_{\gamma_\cardfont{b}})}=
\aleph_{\gamma_\cardfont{b}+1}$. Jika $\cardfont{a}$ bukan penerus dari
sebarang kardinal, maka karena kardinal merupakan ordinal,
$\cardfont{a} = \bigcup_{\cardfont{b} < \cardfont{a}} \cardfont{b} =
\bigcup_{\cardfont{b} < \cardfont{a}}{\aleph_{\gamma_\cardfont{b}}}$,
sehingga $\cardfont{a} = \aleph_\gamma$ dengan
$\gamma = \bigcup_{\cardfont{b} < \cardfont{a}}\gamma_\cardfont{b}$.
\end{proof}

Karena setiap kardinal tak hingga merupakan suatu $\aleph$, hal ini mendorong
kita untuk bertanya: apakah setiap kardinal tak hingga merupakan
suatu~$\beth$? Tentu saja, seandainya \emph{memang} demikian, kardinal-kardinal tak
hingga akan ``berinteraksi secara langsung'' dengan operasi mengambil
himpunan kuasa. Bahkan, kita akan memperoleh pernyataan berikut:

\begin{defish}
\emph{Hipotesis Kontinuum Terampat} (GCH). $\aleph_\alpha = \beth_\alpha$,
untuk semua $\alpha$.
\end{defish}

Lebih jauh, jika GCH berlaku, kita dapat membuat beberapa penyederhanaan
besar pada eksponensiasi kardinal. Secara khusus, kita dapat menunjukkan bahwa
ketika $\cardfont{b} < \cardfont{a}$, nilai
$\cardexpo{\cardfont{a}}{\cardfont{b}}$ dibatasi oleh
$\cardfont{a}\leq \cardexpo{\cardfont{a}}{\cardfont{b}} \leq
\cardsucc{\cardfont{a}}$. Selanjutnya, kita dapat memberikan syarat-syarat
tepat yang menentukan mana dari kedua kemungkinan yang berlaku (yakni,
apakah $\cardfont{a} = \cardexpo{\cardfont{a}}{\cardfont{b}}$ atau
$\cardexpo{\cardfont{a}}{\cardfont{b}} =
\cardsucc{\cardfont{a}}$).\footnote{Syarat tersebut ditentukan oleh
\emph{kofinalitas}.}

Namun, GCH merupakan suatu \emph{hipotesis}, bukan suatu \emph{teorema}.
Pada kenyataannya, \citet{Godel1938} membuktikan bahwa jika $\ZFC$ konsisten,
maka $\ZFC + \text{GCH}$ juga konsisten. Akan tetapi, kemudian terbukti bahwa
kita juga dapat menambahkan $\lnot$GCH pada $\ZFC$. Perhatikan \emph{instans} GCH
tak-sepele yang paling sederhana, yakni:

\begin{defish}
\emph{Hipotesis Kontinuum} (CH). $\aleph_1 = \beth_1$.
\end{defish}

\citet{Cohen1963} membuktikan bahwa jika $\ZFC$ konsisten, maka
$\ZFC + \lnot\text{CH}$ juga konsisten. Jadi, Hipotesis Kontinuum independen
dari $\ZFC$.

Nama Hipotesis Kontinuum digunakan karena ``kontinuum'' merupakan nama lain
bagi garis bilangan real, $\Real$.
\olref[opps]{continuumis2aleph0} memberi tahu kita bahwa
$\card{\Real} = \beth_1$. Jadi, Hipotesis Kontinuum menyatakan bahwa tidak
terdapat suatu kardinal di antara kardinalitas bilangan asli,
$\aleph_0 = \beth_0$, dan kardinalitas kontinuum, $\beth_1$.

Mengingat \emph{independensi} (G)CH dari $\ZFC$, apa yang harus kita katakan
tentang \emph{kebenaran} keduanya? Ada \emph{banyak} hal yang dapat dikatakan.
Memang, banyak penelitian mutakhir yang produktif dalam teori himpunan telah
diarahkan untuk menyelidiki persoalan-persoalan ini. Namun, dua hal singkat
berikut tentu patut ditekankan.

Pertama, hasil-hasil independensi formal ini tidak \emph{langsung}
menyiratkan bahwa nilai kebenaran GCH atau CH \emph{tak tentu}. Mungkin saja
kita hanya perlu menambahkan lebih banyak aksioma yang terasa alami bagi
kita dan yang akan menentukan jawaban atas persoalan tersebut. G\"odel sendiri
mengusulkan bahwa inilah tanggapan yang tepat.

Kedua, independensi CH dari $\ZFC$ tentu \emph{mencolok}, tetapi jelas bukan
sesuatu yang \emph{mustahil dipercaya} (dalam arti harfiah). Pokoknya hanyalah
bahwa, sejauh yang dinyatakan oleh $\ZFC$, berpindah dari kardinal ke
penerusnya mungkin memerlukan alat yang lebih halus daripada sekadar
mengambil himpunan kuasa.
 %The operation of taking powersets moves us from one stage of the
 %hierarchy to its successor stage (from $V_{\alpha}$ to
 %$V_{\alpha+1}$).

Setelah dua pengamatan tersebut, jika Anda ingin mengetahui lebih lanjut,
Anda kini harus beralih kepada berbagai filsuf dan matematikawan yang
mempunyai kepentingan dalam perdebatan ini.\footnote{Namun, Anda mungkin ingin
memulai dengan membaca \citet[\S15.6]{Potter2004}.}

\end{document}
